\documentclass[11pt, a4paper]{article}

\usepackage[utf8]{inputenc}
\usepackage[T1]{fontenc}
\usepackage{amsmath, amssymb, amsthm, mathtools}
\usepackage{bm}
\usepackage{booktabs}
\usepackage{graphicx}
\usepackage{hyperref}
\usepackage[margin=2.5cm]{geometry}
\usepackage{enumitem}
\usepackage{xcolor}
\usepackage{tcolorbox}
\usepackage{float}

% Theorem environments
\theoremstyle{definition}
\newtheorem{definition}{Definition}[section]
\newtheorem{example}{Example}[section]
\theoremstyle{plain}
\newtheorem{theorem}{Theorem}[section]
\newtheorem{proposition}{Proposition}[section]
\newtheorem{corollary}{Corollary}[section]
\theoremstyle{remark}
\newtheorem{remark}{Remark}[section]

% Colored boxes
\newtcolorbox{keyresult}[1][]{
    colback=blue!5!white,
    colframe=blue!60!black,
    fonttitle=\bfseries,
    title={Key Insight},
    #1
}
\newtcolorbox{warningbox}[1][]{
    colback=red!5!white,
    colframe=red!60!black,
    fonttitle=\bfseries,
    title={Warning},
    #1
}
\newtcolorbox{refbox}[1][]{
    colback=gray!5!white,
    colframe=gray!50!black,
    fonttitle=\bfseries,
    title={Go Deeper},
    #1
}
\newtcolorbox{keyidea}[1][]{
    colback=green!5!white,
    colframe=green!50!black,
    fonttitle=\bfseries,
    title={Core Idea},
    #1
}
\newtcolorbox{prereq}[1][]{
    colback=orange!5!white,
    colframe=orange!60!black,
    fonttitle=\bfseries,
    title={Read Before Coding},
    #1
}
\newtcolorbox{objective}[1][]{
    colback=blue!5!white,
    colframe=blue!60!black,
    fonttitle=\bfseries,
    title={What You Will Build},
    #1
}

% Custom commands
\newcommand{\VaR}{\operatorname{VaR}}
\newcommand{\ES}{\operatorname{ES}}
\newcommand{\CVaR}{\operatorname{CVaR}}
\newcommand{\EV}{\mathbb{E}}
\newcommand{\PV}{\mathbb{P}}
\newcommand{\RV}{\mathbb{R}}
\newcommand{\ind}{\boldsymbol{1}}
\newcommand{\PnL}{\operatorname{PnL}}
\newcommand{\CVA}{\operatorname{CVA}}
\newcommand{\EE}{\operatorname{EE}}
\newcommand{\PFE}{\operatorname{PFE}}
\newcommand{\EPE}{\operatorname{EPE}}
\newcommand{\LGD}{\operatorname{LGD}}
\newcommand{\PD}{\operatorname{PD}}

\title{%
    \textbf{CVA and Counterparty Credit Risk} \\[0.5em]
    \large Pricing the Risk That Your Counterparty Defaults
}
\author{Andr\'es Gonz\'alez Ortega}
\date{March 2026}

\begin{document}
\maketitle

\begin{abstract}
An applied treatment of Credit Valuation Adjustment (CVA) and counterparty credit risk ---
the risk that a trading counterparty defaults when it owes you money.
We develop the CVA framework from first principles: what counterparty risk is and why it
differs from market risk and traditional credit risk; the CVA formula as the expected loss
from counterparty default; how to construct exposure profiles via Monte Carlo simulation;
how to extract default probabilities from CDS spreads; the role of netting agreements and
collateral in reducing exposure; the amplifying effect of wrong-way risk; and the regulatory
treatment under Basel~III/IV.
A worked numerical example runs throughout: computing CVA for a 5-year interest rate swap
with quarterly rebalancing, realistic recovery rates, and CDS-implied default probabilities.
We assume familiarity with derivatives pricing, basic credit risk concepts, and Monte Carlo methods.
\end{abstract}

\tableofcontents
\newpage

%% ============================================================================
\section{What Is Counterparty Credit Risk?}
\label{sec:counterparty-risk}
%% ============================================================================

You enter a 5-year interest rate swap with Bank~X\@.
Under this contract, you pay a fixed rate and receive a floating rate on a notional
of \$100 million.
Three years in, Bank~X defaults.
At that moment, the swap has a positive mark-to-market value of \$4 million to you ---
Bank~X owes you money.
But Bank~X is now in bankruptcy.
You will not recover the full \$4 million; perhaps you recover 40\% through the
bankruptcy process, losing \$2.4 million.

This is \textbf{counterparty credit risk}: the risk of loss arising from the default
of a counterparty to a financial contract, when that contract has positive value to you
at the time of default.

\subsection{How It Differs from Other Risks}

Counterparty credit risk sits at the intersection of market risk and credit risk,
but is distinct from both:

\begin{itemize}
    \item \textbf{Market risk}: the risk that the value of a position changes due to
          movements in market variables (interest rates, FX rates, equity prices).
          Market risk exists whether or not anyone defaults.
    \item \textbf{Credit risk} (lending): the risk that a borrower fails to repay a loan.
          The exposure is known and fixed --- the outstanding principal.
    \item \textbf{Counterparty credit risk}: the risk that a trading counterparty defaults
          on a \emph{derivative} contract.
          The critical difference: the exposure is \emph{uncertain} and \emph{two-sided}.
          A swap can have positive or negative value to you depending on how rates move.
          You only lose if the counterparty defaults \emph{and} the contract is in your favour.
\end{itemize}

\begin{keyidea}
Counterparty credit risk has three distinguishing features:
\begin{enumerate}
    \item \textbf{Stochastic exposure}: unlike a loan, the amount at risk changes with
          market conditions.
    \item \textbf{Bilateral nature}: both parties can owe each other money.
          You face counterparty risk only when the mark-to-market is positive to you.
    \item \textbf{Time dimension}: the counterparty can default at any point over the
          life of the trade, and your exposure at that moment is what determines your loss.
\end{enumerate}
\end{keyidea}

\subsection{The Lehman Lesson}

When Lehman Brothers filed for bankruptcy on 15 September 2008, it had approximately
930{,}000 outstanding derivative contracts with counterparties worldwide.
The total notional exceeded \$35 trillion.
Counterparties faced two immediate problems:

\begin{enumerate}
    \item \textbf{Direct losses}: where Lehman owed money on in-the-money contracts,
          counterparties recovered only a fraction (ultimately around 28 cents on the dollar
          for senior unsecured claims).
    \item \textbf{Replacement cost}: counterparties had to re-enter equivalent contracts
          with new counterparties, often at worse market rates because the market had moved
          violently during the crisis.
\end{enumerate}

The total counterparty credit losses from the Lehman default were estimated at
\$5--\$8 billion for its derivative counterparties alone.
This event demonstrated that counterparty risk was not a theoretical concern but a
first-order risk that required explicit pricing and management.

\begin{warningbox}
Before 2008, many banks treated counterparty risk as a secondary concern,
relying on the assumption that large financial institutions would not default.
The Lehman collapse destroyed this assumption permanently.
Since then, CVA has become one of the largest valuation adjustments on bank
trading books, often exceeding the profit margin on individual trades.
\end{warningbox}

\subsection{When Does Counterparty Risk Arise?}

Counterparty credit risk arises in \textbf{over-the-counter (OTC) derivatives} ---
contracts negotiated bilaterally between two parties, as opposed to
exchange-traded instruments where a central clearinghouse guarantees performance.

Common OTC instruments with significant counterparty risk:
\begin{itemize}
    \item Interest rate swaps (the largest market by notional)
    \item Cross-currency swaps
    \item Credit default swaps (CDS)
    \item Equity options and total return swaps
    \item Commodity forwards and swaps
    \item Foreign exchange forwards (beyond short tenors)
\end{itemize}

Since the Dodd--Frank Act (2010) and EMIR in Europe, many standardized OTC derivatives
must be centrally cleared, which transfers counterparty risk to the clearinghouse.
However, a substantial volume of non-standard (``bespoke'') derivatives remains
bilaterally traded, and counterparty risk on these trades must be managed directly.

%% ============================================================================
\section{CVA: The Price of Counterparty Risk}
\label{sec:cva-formula}
%% ============================================================================

Credit Valuation Adjustment (CVA) is the market price of counterparty credit risk.
It quantifies the expected loss due to counterparty default over the life of a trade
or portfolio.

\subsection{The Fundamental CVA Formula}

The unilateral CVA is defined as:
\begin{equation}\label{eq:cva-integral}
    \CVA = (1 - R)\int_0^T \EE(t)\, d\PD(t)
\end{equation}
where:
\begin{itemize}
    \item $R$ is the \textbf{recovery rate} --- the fraction of the exposure recovered
          in the event of default (typically 40\% for senior unsecured claims).
    \item $\EE(t)$ is the \textbf{Expected Exposure} at time $t$ --- the average
          positive value of the contract conditional on surviving to time $t$.
    \item $\PD(t)$ is the \textbf{cumulative default probability} up to time $t$,
          so $d\PD(t)$ represents the probability of default in the infinitesimal interval $[t, t+dt]$.
    \item $T$ is the maturity of the contract.
\end{itemize}

The factor $(1 - R)$ is sometimes called the \textbf{Loss Given Default} ($\LGD$).
If the counterparty defaults and you have \$1 of exposure, you lose \$(1 - R).

\subsection{Discrete Approximation}

In practice, we partition $[0, T]$ into $N$ time intervals with dates
$t_0 = 0, t_1, \ldots, t_N = T$, and approximate:
\begin{equation}\label{eq:cva-discrete}
    \CVA = (1 - R)\sum_{i=1}^{N} \EE(t_i)\cdot\bigl[\PD(t_i) - \PD(t_{i-1})\bigr]
\end{equation}
Each term in the sum is the expected loss from a default occurring in the interval
$(t_{i-1}, t_i]$: the exposure at that time, multiplied by the probability of default
in that interval, multiplied by the loss given default.

\begin{keyresult}
CVA is the \textbf{risk-neutral expected loss} from counterparty default.
It is subtracted from the risk-free (counterparty-risk-free) price of the contract.
If the risk-free price of a swap is \$1{,}000{,}000 and CVA is \$50{,}000,
then the fair value accounting for counterparty risk is \$950{,}000.
\[
    V_{\text{risky}} = V_{\text{risk-free}} - \CVA
\]
The counterparty effectively ``charges'' you for the credit risk they impose on the trade.
\end{keyresult}

\subsection{Worked Example: CVA for an Interest Rate Swap}

\begin{example}[CVA calculation for a 5-year swap]\label{ex:cva-swap}
Consider a 5-year pay-fixed, receive-floating interest rate swap with notional \$100M\@.
Suppose we have computed the following exposure profile and default probabilities
(from CDS spreads, as derived in Section~\ref{sec:default-probability}):

\begin{center}
\begin{tabular}{@{}cccc@{}}
\toprule
\textbf{Year} $t_i$ & $\EE(t_i)$ (\$M) & $\PD(t_i)$ & $\PD(t_i) - \PD(t_{i-1})$ \\
\midrule
1 & 1.2 & 1.65\% & 1.65\% \\
2 & 2.1 & 3.25\% & 1.60\% \\
3 & 2.8 & 4.80\% & 1.55\% \\
4 & 2.3 & 6.30\% & 1.50\% \\
5 & 1.5 & 7.75\% & 1.45\% \\
\bottomrule
\end{tabular}
\end{center}

With recovery rate $R = 40\%$:
\begin{align*}
    \CVA &= (1 - 0.40)\bigl[1.2 \times 0.0165 + 2.1 \times 0.0160 + 2.8 \times 0.0155 \\
         &\quad + 2.3 \times 0.0150 + 1.5 \times 0.0145\bigr] \\
         &= 0.60 \times [0.0198 + 0.0336 + 0.0434 + 0.0345 + 0.0218] \\
         &= 0.60 \times 0.1531 \\
         &= 0.0919 \text{ (\$M)}
\end{align*}
The CVA is approximately \$92{,}000 on a \$100M notional swap.
This amount would be charged to the counterparty at trade inception as a credit
adjustment, or reserved on the trading book as a valuation adjustment.
\end{example}

\begin{remark}
The exposure profile has a characteristic ``hump'' shape for interest rate swaps:
exposure starts near zero (the swap is at-the-money at inception),
rises as rates drift and dispersion grows, then declines as the remaining
cash flows decrease toward maturity.
The peak exposure typically occurs at approximately one-third to one-half of the
swap's tenor.
\end{remark}

\subsection{Bilateral CVA and DVA}

The formula above is \emph{unilateral} CVA --- it accounts for the counterparty's
default only.
But you can also default on your counterparty.
The \textbf{Debit Valuation Adjustment} (DVA) is the symmetric quantity:
\[
    \text{DVA} = (1 - R_{\text{own}})\sum_{i=1}^{N} \text{NEE}(t_i)\cdot
    \bigl[\PD_{\text{own}}(t_i) - \PD_{\text{own}}(t_{i-1})\bigr]
\]
where $\text{NEE}(t)$ is the negative expected exposure (what you owe the counterparty)
and $\PD_{\text{own}}$ is your own default probability.

The bilateral adjustment is:
\[
    V_{\text{risky}} = V_{\text{risk-free}} - \CVA + \text{DVA}
\]

\begin{warningbox}
DVA is controversial: it implies that your derivative position \emph{increases} in value
when your own credit quality deteriorates.
If your CDS spread widens, DVA rises, producing a ``profit.''
This is economically logical (your counterparty benefits from your higher default probability)
but practically problematic --- you cannot monetize your own default.
Under IFRS~13 and ASC~820, DVA must be recognized in fair value, but Basel~III excludes
DVA from regulatory capital.
\end{warningbox}

%% ============================================================================
\section{Exposure Profiles}
\label{sec:exposure-profiles}
%% ============================================================================

The exposure profile $\EE(t)$ is the engine of CVA computation.
To compute it, we need to simulate the future value of the contract under
many scenarios.

\subsection{Key Exposure Metrics}

Let $V(t)$ denote the mark-to-market value of the contract at time $t$.
Your credit exposure is $\max(V(t), 0)$ --- you are exposed only when
the counterparty owes you money.

\begin{definition}[Exposure Metrics]\label{def:exposure-metrics}
The key exposure measures are:
\begin{align}
    \EE(t) &= \EV\bigl[\max(V(t), 0)\bigr]
    && \text{(Expected Exposure)} \label{eq:ee}\\[4pt]
    \PFE_\alpha(t) &= \inf\bigl\{x : \PV\bigl[\max(V(t), 0) \leq x\bigr] \geq \alpha\bigr\}
    && \text{(Potential Future Exposure)} \label{eq:pfe}\\[4pt]
    \EPE &= \frac{1}{T}\int_0^T \EE(t)\, dt
    && \text{(Expected Positive Exposure)} \label{eq:epe}
\end{align}
\end{definition}

\begin{itemize}
    \item $\EE(t)$ is the \textbf{average positive value} across scenarios at time $t$.
          This is the key input to the CVA formula.
    \item $\PFE_\alpha(t)$ is the $\alpha$-quantile of the positive exposure distribution.
          With $\alpha = 97.5\%$, it answers: ``what is the worst-case exposure exceeded
          only 2.5\% of the time?''
          PFE is used for \textbf{credit limit management} --- setting maximum exposure
          to any single counterparty.
    \item $\EPE$ is the time-averaged expected exposure.
          It is used in the \textbf{Basel regulatory capital} calculation as a simplified
          measure of overall exposure.
\end{itemize}

\subsection{Monte Carlo Exposure Simulation}

The standard approach to computing exposure profiles is Monte Carlo simulation:

\begin{enumerate}
    \item \textbf{Simulate risk factors}: generate $M$ scenarios (e.g., $M = 10{,}000$)
          of the underlying risk factors (interest rates, FX rates, etc.)\ at each
          future date $t_1, \ldots, t_N$.
    \item \textbf{Value the contract}: for each scenario $j$ and each date $t_i$,
          compute the mark-to-market value $V^{(j)}(t_i)$ of the contract.
    \item \textbf{Compute positive exposure}: set
          $E^{(j)}(t_i) = \max\bigl(V^{(j)}(t_i), 0\bigr)$.
    \item \textbf{Aggregate}:
          \[
              \EE(t_i) = \frac{1}{M}\sum_{j=1}^{M} E^{(j)}(t_i), \qquad
              \PFE_{97.5\%}(t_i) = \text{97.5th percentile of } \{E^{(j)}(t_i)\}_{j=1}^{M}
          \]
\end{enumerate}

\begin{example}[Exposure simulation for a 5-year swap]\label{ex:exposure-sim}
Consider a 5-year pay-fixed, receive-floating swap on 3-month SOFR, notional \$100M,
fixed rate 3.5\%, quarterly payment dates.

\textbf{Simulation}: generate 10{,}000 interest rate paths using a Hull--White model
calibrated to the current yield curve and swaption volatilities.
At each quarterly date (20 dates total), price the residual swap under each scenario.

\textbf{Representative results}:
\begin{center}
\begin{tabular}{@{}cccc@{}}
\toprule
\textbf{Year} & $\EE(t)$ (\$M) & $\PFE_{97.5\%}(t)$ (\$M) & $\PFE_{97.5\%}/\EE$ ratio \\
\midrule
0.5  & 0.7  & 2.1  & 3.0 \\
1.0  & 1.2  & 3.8  & 3.2 \\
2.0  & 2.1  & 6.5  & 3.1 \\
3.0  & 2.8  & 8.2  & 2.9 \\
4.0  & 2.3  & 6.1  & 2.7 \\
5.0  & 1.5  & 3.9  & 2.6 \\
\bottomrule
\end{tabular}
\end{center}

\textbf{Observations}:
\begin{enumerate}
    \item The exposure profile has a \textbf{hump shape}: it rises as rate uncertainty
          grows, peaks around year 2--3, and declines as the remaining cash flows
          shrink (``pull-to-par'' effect).
    \item PFE is roughly 2.5--3.2 times EE\@.
          The ratio is higher at shorter horizons where the distribution of exposure
          is more skewed (many paths still near zero, but some have moved significantly).
    \item At maturity ($t = 5$), both EE and PFE drop --- only the final coupon
          exchange remains.
\end{enumerate}
\end{example}

\begin{keyidea}
The exposure profile depends heavily on the \textbf{type of derivative}:
\begin{itemize}
    \item \textbf{Interest rate swaps}: hump-shaped profile (diffusion then pull-to-par).
    \item \textbf{Cross-currency swaps}: exposure grows toward maturity because of
          the final notional exchange, producing an \emph{increasing} profile.
    \item \textbf{FX forwards}: exposure grows monotonically until settlement.
    \item \textbf{Options}: exposure is bounded by the option value;
          for purchased options, it can only decrease after the initial premium is paid.
\end{itemize}
Understanding these shapes is essential for intuition about which trades
drive CVA in a portfolio.
\end{keyidea}

%% ============================================================================
\section{Default Probability from CDS Spreads}
\label{sec:default-probability}
%% ============================================================================

The CVA formula requires the counterparty's default probability curve $\PD(t)$.
The market-standard source is \textbf{Credit Default Swap (CDS) spreads},
which embed the market's view of default risk.

\subsection{CDS Basics}

A CDS is an insurance contract on a reference entity's default.
The protection buyer pays a periodic premium (the CDS spread $s$) and,
if the reference entity defaults, receives $(1 - R)$ times the notional.
In a fair CDS, the expected premium payments equal the expected default payment:
\[
    s \sum_{i=1}^{N} \Delta t_i \cdot Q(t_i)
    \approx (1 - R)\sum_{i=1}^{N} \bigl[Q(t_{i-1}) - Q(t_i)\bigr]
\]
where $Q(t) = 1 - \PD(t)$ is the survival probability and we have omitted discounting
for clarity.

\subsection{Hazard Rate and Survival Probability}

Under a constant hazard rate model, the default intensity is a constant $\lambda$:
\begin{align}
    Q(t) &= e^{-\lambda t} && \text{(Survival probability)} \label{eq:survival}\\
    \PD(t) &= 1 - e^{-\lambda t} && \text{(Cumulative default probability)} \label{eq:pd}
\end{align}

The probability of default in a small interval $[t, t + dt]$ given survival to $t$
is $\lambda\, dt$.

For a flat CDS spread $s$ and constant recovery $R$, the implied hazard rate is
approximately:
\begin{equation}\label{eq:hazard-approx}
    \lambda \approx \frac{s}{1 - R}
\end{equation}

This approximation is excellent for investment-grade names (where $\lambda$ is small)
and provides the key intuition: the hazard rate is the CDS spread divided by the
loss given default.

\begin{example}[Default probability from CDS spread]\label{ex:cds-bootstrap}
A counterparty has a 5-year CDS spread of 100 basis points (bps) and
assumed recovery rate $R = 40\%$.

\textbf{Step 1}: compute the hazard rate:
\[
    \lambda = \frac{s}{1 - R} = \frac{0.0100}{1 - 0.40} = \frac{0.0100}{0.60} = 0.01\overline{6} = 1.67\% \text{ per year}
\]

\textbf{Step 2}: compute survival and default probabilities:
\begin{center}
\begin{tabular}{@{}cccc@{}}
\toprule
\textbf{Year} $t$ & $Q(t) = e^{-\lambda t}$ & $\PD(t) = 1 - Q(t)$ & $\PD(t_i) - \PD(t_{i-1})$ \\
\midrule
1 & 98.35\% & 1.65\% & 1.65\% \\
2 & 96.73\% & 3.27\% & 1.62\% \\
3 & 95.12\% & 4.88\% & 1.61\% \\
4 & 93.54\% & 6.46\% & 1.58\% \\
5 & 91.98\% & 8.02\% & 1.56\% \\
\bottomrule
\end{tabular}
\end{center}

The 5-year cumulative default probability is approximately 8.0\%.
The marginal (annual) default probability declines slightly because the hazard rate
is applied to a shrinking pool of surviving paths.
\end{example}

\subsection{Bootstrapping from the CDS Term Structure}

In practice, CDS spreads are quoted at multiple tenors (1Y, 3Y, 5Y, 7Y, 10Y),
and the hazard rate is \emph{not} constant.
The bootstrap procedure extracts a piecewise-constant hazard rate curve:

\begin{enumerate}
    \item Start with the 1Y CDS spread $s_1$.
          Solve for $\lambda_1$ such that the 1Y CDS is fairly priced.
    \item Given $\lambda_1$, use the 3Y CDS spread $s_3$ to solve for $\lambda_2$
          (the hazard rate for the interval $(1, 3]$).
    \item Continue for each subsequent tenor, solving for the hazard rate in each interval.
\end{enumerate}

\begin{example}[Bootstrap from CDS term structure]\label{ex:cds-term}
Suppose the following CDS spreads (with $R = 40\%$):

\begin{center}
\begin{tabular}{@{}cc@{\qquad}cc@{}}
\toprule
\textbf{Tenor} & \textbf{CDS spread (bps)} & \textbf{Implied } $\lambda$ & $\PD$ \text{(cumulative)} \\
\midrule
1Y  &  60  & 1.00\% & 1.00\% \\
3Y  &  80  & 1.43\% & 3.79\% \\
5Y  & 100  & 1.82\% & 7.24\% \\
7Y  & 110  & 1.50\% & 9.50\% \\
10Y & 120  & 1.67\% & 14.2\% \\
\bottomrule
\end{tabular}
\end{center}

The increasing CDS term structure implies that the market expects the hazard rate
to be higher at longer horizons --- a normal credit curve.
An inverted curve (decreasing spreads) would suggest near-term distress.
\end{example}

\begin{keyresult}
CDS spreads are the market's most liquid and direct measure of default risk.
The hazard rate $\lambda \approx s / (1 - R)$ translates spreads into default
probabilities.
These probabilities are \textbf{risk-neutral} (they embed a credit risk premium),
so they are typically higher than historical default rates.
For CVA pricing purposes, we use risk-neutral probabilities --- the same principle
as using risk-neutral pricing for derivatives.
\end{keyresult}

%% ============================================================================
\section{Netting and Collateral}
\label{sec:netting-collateral}
%% ============================================================================

In practice, CVA is computed not on individual trades but on \textbf{netting sets} ---
groups of trades with the same counterparty governed by a single legal agreement.
Netting and collateral are the two most powerful tools for reducing counterparty exposure.

\subsection{Netting Agreements}

An ISDA Master Agreement with a close-out netting clause establishes that, upon default,
all trades under the agreement are terminated and their values are netted to a
single amount.
If you have 50 trades with a counterparty, some with positive and some with negative
mark-to-market, your exposure is:
\[
    E_{\text{net}}(t) = \max\!\left(\sum_{k=1}^{50} V_k(t),\; 0\right)
\]
\textbf{not} the sum of individual positive exposures:
\[
    E_{\text{gross}}(t) = \sum_{k=1}^{50} \max\bigl(V_k(t), 0\bigr)
\]

The netting benefit is:
\[
    E_{\text{gross}}(t) - E_{\text{net}}(t) \geq 0
\]
which is always non-negative (by the subadditivity of the max function).

\begin{keyidea}
Netting can dramatically reduce exposure.
Consider a counterparty where you have an interest rate swap worth $+\$5$M
and a cross-currency swap worth $-\$3$M\@.
\begin{itemize}
    \item \textbf{Without netting}: exposure = $\$5$M (the negative swap is ignored).
    \item \textbf{With netting}: exposure = $\max(5 - 3, 0) = \$2$M\@.
\end{itemize}
A 60\% reduction.
For large portfolios with many offsetting positions, the netting benefit
can exceed 80\% of the gross exposure.
\end{keyidea}

\begin{example}[Netting benefit]\label{ex:netting}
A bank has the following five trades with counterparty~C, all under a single
ISDA Master Agreement:

\begin{center}
\begin{tabular}{@{}lc@{}}
\toprule
\textbf{Trade} & \textbf{MtM (\$M)} \\
\midrule
5Y IRS (pay fixed)        & $+3.2$ \\
3Y IRS (receive fixed)    & $-1.8$ \\
7Y cross-currency swap    & $+4.5$ \\
FX forward (6 months)     & $-0.6$ \\
Equity option (purchased) & $+1.1$ \\
\bottomrule
\end{tabular}
\end{center}

\textbf{Gross exposure}: $3.2 + 0 + 4.5 + 0 + 1.1 = \$8.8$M (sum of positive values).

\textbf{Net exposure}: $\max(3.2 - 1.8 + 4.5 - 0.6 + 1.1,\; 0) = \max(6.4,\; 0) = \$6.4$M.

\textbf{Netting benefit}: $\$8.8$M $- \$6.4$M $= \$2.4$M, a 27\% reduction.
The benefit would be larger if the portfolio had more balanced positive and negative values.
\end{example}

\subsection{Collateral (CSA Agreements)}

A Credit Support Annex (CSA) to the ISDA Master Agreement requires counterparties
to post collateral (cash or high-quality securities) when the mark-to-market
exceeds a specified threshold.

Key CSA parameters:
\begin{itemize}
    \item \textbf{Threshold} ($H$): exposure below this level requires no collateral.
    \item \textbf{Minimum Transfer Amount} (MTA): collateral calls below this amount
          are not made (to avoid operational costs on small amounts).
    \item \textbf{Margin frequency}: how often collateral is recalculated (daily, weekly).
    \item \textbf{Independent amount}: an initial margin posted regardless of exposure
          (similar to exchange-traded initial margin).
\end{itemize}

With a CSA in place, the \textbf{collateralized exposure} at time $t$ is approximately:
\begin{equation}\label{eq:collateral-exposure}
    E_{\text{coll}}(t) = \max\bigl(V(t) - C(t),\; 0\bigr)
\end{equation}
where $C(t)$ is the collateral held at time $t$.

\subsection{Margin Period of Risk}

Even with daily margining, collateral is not instantaneous.
The \textbf{Margin Period of Risk} (MPoR) is the time between the last successful
collateral exchange and the close-out of positions following a default.
During this period, the exposure can increase without being covered by collateral.

\begin{warningbox}
The Basel Committee sets a minimum MPoR of 10 business days for bilateral trades
and 5 business days for centrally cleared trades.
During these 10 days, the counterparty may default, and the exposure can grow significantly.
The CVA on a collateralized trade is therefore \emph{not} zero --- it reflects the
exposure that can build up during the MPoR\@.
For volatile products or stressed markets, the 10-day MPoR exposure can still be
substantial.
\end{warningbox}

\begin{keyresult}
The combined effect of netting and collateral reduces CVA dramatically:
\begin{center}
\begin{tabular}{@{}lcc@{}}
\toprule
\textbf{Configuration} & \textbf{EPE (\$M)} & \textbf{CVA (\$K)} \\
\midrule
Gross (no netting)           & 8.8 & 440 \\
Netted (ISDA)                & 4.2 & 210 \\
Netted + CSA ($H = 0$, daily) & 0.8 &  40 \\
\bottomrule
\end{tabular}
\end{center}
Representative values for a diversified portfolio.
Netting reduces CVA by ${\sim}50\%$; adding collateral can reduce it by a further
${\sim}80\%$.
\end{keyresult}

%% ============================================================================
\section{Wrong-Way Risk}
\label{sec:wrong-way-risk}
%% ============================================================================

The standard CVA formula~\eqref{eq:cva-integral} assumes that the counterparty's
default probability is \textbf{independent} of the exposure.
This assumption is often violated, sometimes catastrophically.

\subsection{Definition}

\begin{definition}[Wrong-Way Risk]\label{def:wwr}
\textbf{Wrong-way risk} (WWR) arises when the counterparty's default probability
is positively correlated with the exposure to that counterparty:
\[
    \operatorname{Corr}\bigl(\max(V(t), 0),\; \ind_{\{\tau \leq t\}}\bigr) > 0
\]
where $\tau$ is the default time.
In other words, the counterparty is \emph{more likely to default precisely when
it owes you the most}.
\end{definition}

The opposite --- \textbf{right-way risk} --- occurs when default probability is
negatively correlated with exposure (the counterparty defaults when it owes you less).

\subsection{Examples of Wrong-Way Risk}

\textbf{Example 1: Protection from a correlated bank.}
You buy credit protection (CDS) from Bank~A on a mortgage-backed security (MBS).
If the housing market crashes:
\begin{itemize}
    \item The MBS value drops, making your CDS more valuable (exposure rises).
    \item Bank~A, heavily exposed to mortgages, becomes more likely to default.
\end{itemize}
This is exactly what happened with AIG in 2008: AIG sold CDS protection on
mortgage-related CDOs, and its creditworthiness deteriorated simultaneously
with the rise in claims against it.

\textbf{Example 2: FX swap with an emerging-market bank.}
You enter a cross-currency swap with a Brazilian bank where you receive BRL
and pay USD\@.
If the BRL depreciates sharply:
\begin{itemize}
    \item The swap value to you increases (you receive more BRL per USD).
    \item The Brazilian bank faces currency stress, and its default probability rises.
\end{itemize}

\textbf{Example 3: Equity swap with a corporate.}
You enter a total return swap on Company~C's stock with Company~C as counterparty.
If the stock price drops, Company~C's creditworthiness declines, but the swap
has positive value to you (you are short the stock via the swap).

\subsection{Impact on CVA}

Under wrong-way risk, the true CVA is:
\begin{equation}\label{eq:cva-wwr}
    \CVA_{\text{WWR}} = (1 - R)\int_0^T \EV\!\bigl[\max(V(t), 0) \mid \tau = t\bigr]\, d\PD(t)
\end{equation}
Note the key difference from~\eqref{eq:cva-integral}: the expectation is now
\emph{conditional on default at time $t$}, not unconditional.
Under WWR, the conditional exposure $\EV[\max(V(t), 0) \mid \tau = t]$
is higher than the unconditional $\EE(t)$, so:
\[
    \CVA_{\text{WWR}} > \CVA_{\text{independent}}
\]

\begin{warningbox}
Wrong-way risk can increase CVA by \textbf{multiples}, not merely percentages.
In extreme cases (e.g., buying protection from a counterparty on a highly correlated
reference), the CVA under WWR can be 2--5 times the value computed under independence.
The standard CVA formula without WWR adjustment can be dangerously misleading.
\end{warningbox}

\subsection{Modelling Approaches}

Several approaches exist to incorporate wrong-way risk:

\begin{enumerate}
    \item \textbf{Alpha multiplier}: a regulatory shortcut.
          Multiply the independent CVA by a factor $\alpha > 1$ (Basel sets $\alpha = 1.4$
          as the default, though banks with IMM approval may use their own estimates).
    \item \textbf{Copula models}: model the joint distribution of the exposure and
          default time using a copula (e.g., Gaussian copula with a correlation parameter).
          Simple to implement, but the correlation parameter is hard to calibrate.
    \item \textbf{Structural models}: link the counterparty's default to the same
          risk factors driving the exposure (e.g., interest rates, FX rates, equity prices).
          The counterparty defaults when its asset value hits a barrier,
          and the same market moves that increase exposure also push the counterparty
          toward the barrier.
    \item \textbf{Jump-to-default}: model a sudden jump in market variables upon default
          (e.g., the counterparty's bonds drop to recovery value, correlated assets move).
\end{enumerate}

\begin{keyresult}
Wrong-way risk is the ``hidden amplifier'' of counterparty risk.
The standard CVA formula assumes independence between exposure and default,
but many real trades exhibit positive correlation.
Ignoring WWR leads to systematic underestimation of CVA\@.
The Lehman default and AIG bailout both demonstrated that wrong-way risk
materializes precisely in systemic crises when it matters most.
\end{keyresult}

%% ============================================================================
\section{Regulatory CVA (Basel III/IV)}
\label{sec:regulatory-cva}
%% ============================================================================

Following the 2008 financial crisis, regulators recognized that CVA losses ---
not just outright counterparty defaults --- caused enormous losses for banks.
Two-thirds of the counterparty credit losses during the crisis came from
mark-to-market CVA losses (changes in CVA due to spread widening),
not from actual defaults.

Basel~III introduced a dedicated \textbf{CVA capital charge} to capture this risk.
Basel~IV (the ``finalized Basel~III'' framework, effective from January 2025)
revised the CVA capital framework significantly.

\subsection{The CVA Capital Charge}

The CVA capital charge is separate from the \textbf{counterparty default risk charge}
(which covers actual losses from default, computed under CCR rules).
The CVA charge specifically covers the risk of \textbf{mark-to-market losses}
due to changes in the CVA --- that is, changes in credit spreads, exposure,
and their interaction.

\subsection{SA-CVA: Standardized Approach}

The Standardized Approach for CVA (SA-CVA) is a sensitivity-based method,
structurally similar to the FRTB Standardized Approach for market risk.

\textbf{Key features}:
\begin{enumerate}
    \item Compute sensitivities of the aggregate CVA to risk factors:
          counterparty credit spreads, interest rates, FX rates, equity prices,
          and commodity prices.
    \item Apply prescribed risk weights to each sensitivity.
    \item Aggregate using a correlation structure specified by the regulator
          to produce the total capital charge.
\end{enumerate}

The SA-CVA capital for a given risk class is computed in two components:
\begin{align}
    K_b &= \sqrt{\sum_k W_k^2 + \sum_{k}\sum_{l \neq k} \rho_{kl}\, W_k\, W_l}
    && \text{(within-bucket)} \label{eq:sa-cva-bucket}\\
    K &= \sqrt{\sum_b K_b^2 + \sum_b \sum_{c \neq b} \gamma_{bc}\, K_b\, K_c}
    && \text{(across-bucket)} \label{eq:sa-cva-total}
\end{align}
where $W_k$ are the weighted sensitivities, $\rho_{kl}$ are within-bucket correlations,
and $\gamma_{bc}$ are across-bucket correlations.

\subsection{BA-CVA: Basic Approach}

For banks that cannot compute CVA sensitivities, the Basic Approach (BA-CVA)
provides a simplified formula.
The reduced version is:
\begin{equation}\label{eq:ba-cva}
    K_{\text{BA-CVA}} = \rho \cdot \sum_c \text{SC}_c
    + \sqrt{1 - \rho^2}\cdot\sqrt{\sum_c \text{SC}_c^2}
\end{equation}
where:
\begin{itemize}
    \item $\text{SC}_c$ is the \textbf{standalone capital} for counterparty $c$,
          a function of the EAD, maturity, counterparty risk weight, and supervisory
          discount factor.
    \item $\rho = 50\%$ is the supervisory correlation, reflecting the systematic
          component of CVA risk.
    \item The first term captures systematic CVA risk (correlated spread movements);
          the second captures idiosyncratic risk (counterparty-specific spread movements).
\end{itemize}

The standalone capital for each counterparty is:
\[
    \text{SC}_c = \frac{1}{\alpha}\cdot \text{RW}_c \cdot M_c \cdot \text{EAD}_c \cdot \text{DF}_c
\]
where $\text{RW}_c$ is the supervisory risk weight (based on credit rating and sector),
$M_c$ is the effective maturity, $\text{EAD}_c$ is the exposure at default,
$\text{DF}_c$ is a supervisory discount factor, and $\alpha = 1.4$.

\begin{refbox}
\begin{itemize}
    \item Basel Committee on Banking Supervision (2020).
          ``Minimum capital requirements for CVA risk.'' MAR50 in the consolidated framework.
          Defines SA-CVA and BA-CVA in full detail.
    \item For the interaction with FRTB market risk capital, see MAR10--MAR33.
\end{itemize}
\end{refbox}

\subsection{Comparison of Approaches}

\begin{center}
\begin{tabular}{@{}p{3.5cm}p{4.8cm}p{4.8cm}@{}}
\toprule
& \textbf{SA-CVA} & \textbf{BA-CVA} \\
\midrule
\textbf{Complexity}     & High: requires CVA sensitivities to all risk factors
                        & Low: uses EAD, maturity, and supervisory risk weights \\[4pt]
\textbf{Risk sensitivity} & High: captures spread, IR, FX, equity risk
                          & Moderate: captures spread risk only \\[4pt]
\textbf{Hedging recognition} & Recognizes eligible CVA hedges (CDS, index CDS)
                              & Does not recognize hedges \\[4pt]
\textbf{Capital level}  & Generally lower (more risk-sensitive)
                        & Generally higher (conservative) \\[4pt]
\textbf{Prerequisites}  & Ability to compute CVA and its sensitivities
                        & Minimal infrastructure \\
\bottomrule
\end{tabular}
\end{center}

\begin{keyresult}
The regulatory CVA framework has two important consequences for banks:
\begin{enumerate}
    \item \textbf{Capital cost}: the CVA capital charge is a significant component
          of total trading book capital, often 10--30\% of the total market risk charge.
          This cost is passed on to clients through wider spreads on OTC derivatives.
    \item \textbf{Incentive to mitigate}: the framework incentivizes netting,
          collateralization, and central clearing, all of which reduce the CVA charge.
          Banks with SA-CVA approval also benefit from hedging recognition,
          creating an incentive to actively hedge CVA risk with CDS.
\end{enumerate}
\end{keyresult}

%% ============================================================================
\section{References}
\label{sec:references}
%% ============================================================================

The following texts provide the mathematical foundations and practical guidance
for CVA and counterparty credit risk.

\begin{enumerate}[label={[\arabic*]}, leftmargin=2cm]
    \item \textbf{Gregory, J.}\ (2020).
          \emph{The xVA Challenge: Counterparty Credit Risk, Funding, Collateral, Capital and Initial Margin.}
          4th edition, Wiley.
          The definitive practitioner reference covering CVA, DVA, FVA, KVA,
          and all valuation adjustments.
          Chapters 8--12 cover exposure simulation, CVA computation, and wrong-way risk in detail.

    \item \textbf{Pykhtin, M.\ \& Zhu, S.H.}\ (2007).
          ``A Guide to Modelling Counterparty Credit Risk.''
          \emph{GARP Risk Review}, July/August 2007.
          An accessible introduction to exposure modelling, netting,
          and the computation of EE, PFE, and EPE\@.

    \item \textbf{Brigo, D.\ \& Morini, M.}\ (2005).
          ``Dangers of Bilateral Counterparty Risk: the fundamental impact
          of closeout conventions.''
          Working paper.
          Explores the subtleties of bilateral CVA, DVA, and closeout
          assumptions.

    \item \textbf{Brigo, D., Morini, M.\ \& Pallavicini, A.}\ (2013).
          \emph{Counterparty Credit Risk, Collateral and Funding:
          With Pricing Cases for All Asset Classes.}
          Wiley.
          Rigorous mathematical treatment of CVA with applications
          to interest rate, FX, credit, and equity derivatives.

    \item \textbf{Basel Committee on Banking Supervision}\ (2020).
          ``Minimum capital requirements for CVA risk (MAR50).''
          \url{https://www.bis.org/basel_framework/}.
          The official regulatory standard for SA-CVA and BA-CVA capital charges.

    \item \textbf{Green, A.}\ (2015).
          \emph{XVA: Credit, Funding and Capital Valuation Adjustments.}
          Wiley.
          Practical guide to implementing CVA and related adjustments,
          with worked examples and code.

    \item \textbf{Cesari, G., Aquilina, J., Charpillon, N., Filipovi\'c, Z.,
          Lee, G.\ \& Manda, I.}\ (2009).
          \emph{Modelling, Pricing, and Hedging Counterparty Credit Exposure.}
          Springer.
          Detailed technical treatment of exposure computation, simulation methods,
          and computational challenges.

    \item \textbf{Hull, J.\ \& White, A.}\ (2012).
          ``CVA and Wrong-Way Risk.''
          \emph{Financial Analysts Journal}, 68(5), 58--69.
          Clear exposition of wrong-way risk modelling with practical examples.
\end{enumerate}

\begin{prereq}
Before implementing CVA in code, ensure you are comfortable with:
\begin{itemize}
    \item Interest rate swap pricing and the mechanics of cash flow exchanges.
    \item Monte Carlo simulation of interest rate and FX paths
          (Hull--White, LMM, or similar models).
    \item CDS pricing and the bootstrap of default probability curves.
    \item Netting and collateral mechanics (ISDA Master Agreement, CSA).
    \item Basic credit risk concepts: default probability, recovery rate, loss given default.
\end{itemize}
\end{prereq}

\vfill
\hrule
\smallskip
\noindent\textit{Andr\'es Gonz\'alez Ortega --- Risk Analyst Project} \hfill \textit{March 2026}

\end{document}
